% Source of the paper since #963. `paper.md` is frozen and will be removed once
% the LaTeX version is approved.
%
% Build from the root of the repo:
% > make -C papers/Optimal_strategy_for_racket_sports latex
\documentclass[11pt, reqno]{amsart}
\usepackage{amsfonts, amssymb, amscd}
%\usepackage{amsrefs}
\usepackage{graphicx}
\usepackage{hyperref}
\usepackage{slashed}
\usepackage{fullpage}
% Prevent table repositioning.
\usepackage{float}
% For textcolor.
\usepackage{xcolor}
% For blackboard bold `1`.
\usepackage{bbold}

\usepackage{booktabs}
\usepackage{tabularx}
\usepackage{longtable, array}

\usepackage{algorithm}
\usepackage{algpseudocode}

\usepackage{makecell}

%\usepackage{titlesec}

%\titleformat*{\section}{\LARGE\bfseries}
%\titleformat*{\subsection}{\Large\bfseries}
%\titleformat*{\subsubsection}{\large\bfseries}
%\titleformat*{\paragraph}{\large\bfseries}
%\titleformat*{\subparagraph}{\large\bfseries}

% If you need math, theorems, etc.
\usepackage{amsmath, amssymb, amsthm}
\usepackage{tikz}
\usetikzlibrary{
  positioning,
  calc,
  arrows.meta,
  decorations.pathreplacing
}

\usepackage[useregional]{datetime2}
\DTMlangsetup[en-US]{zone=eastern,mapzone}

% --- Theorem environments (acmart loads amsthm; this is safe) ---
\theoremstyle{plain}
\newtheorem{theorem}{Theorem}[section]
\newtheorem{lemma}[theorem]{Lemma}
\newtheorem{proposition}[theorem]{Proposition}
\newtheorem{corollary}[theorem]{Corollary}

\theoremstyle{definition}
\newtheorem{definition}[theorem]{Definition}

\theoremstyle{remark}
\newtheorem{remark}[theorem]{Remark}

\newtheorem{assumption}{Assumption}[section]

\input{./helpers_root/dev_scripts_helpers/documentation/latex_abbrevs.sty}

% Python style for highlighting
\usepackage{listings}

\newcommand{\pythonstyle}{\lstset{ language=Python, basicstyle=\ttm, morekeywords={self}, % Add keywords here
keywordstyle=\ttb
\color{deepblue}
, emph={MyClass,__init__}, % Custom highlighting
emphstyle=\ttb
\color{deepred}
, % Custom highlighting style
stringstyle=
\color{deepgreen}
, frame=tb, % Any extra options here
showstringspaces=false }}

% Python environment.
\lstnewenvironment{python}[1][]{ \pythonstyle \lstset{#1} }{}

% Python for external files.
\newcommand{\pythonexternal}[2][]{{ \pythonstyle \lstinputlisting[#1]{#2}}}

% Python for inline.
\newcommand{\pythoninline}[1]{{\pythonstyle\lstinline!#1!}}

\usepackage{listings}
\usepackage{xcolor}

\definecolor{codegreen}{rgb}{0,0.6,0}
\definecolor{codegray}{rgb}{0.5,0.5,0.5}
\definecolor{codepurple}{rgb}{0.58,0,0.82}
\definecolor{backcolour}{rgb}{0.95,0.95,0.92}

\lstdefinestyle{mystyle}{ backgroundcolor=
\color{backcolour}
, commentstyle=
\color{codegreen}
, keywordstyle=
\color{magenta}
, numberstyle=\tiny
\color{codegray}
, stringstyle=
\color{codepurple}
, basicstyle=\ttfamily\footnotesize, breakatwhitespace=false, breaklines=true, captionpos=b,
keepspaces=true, numbers=left, numbersep=5pt, showspaces=false, showstringspaces=false,
showtabs=false, tabsize=2 }

\lstset{style=mystyle}

% end python

%

\providecommand{\tightlist}{%
\setlength{\itemsep}{0pt}
\setlength{\parskip}{0pt}}

% Use this if using `\contrib[]{...}`.
%\makeatletter\let\@wraptoccontribs\wraptoccontribs\makeatother
% https://tex.stackexchange.com/questions/418547/equal-contribution-using-thanks-with-llncs-class#418563
\makeatletter
\newcommand{\printfnsymbol}[1]{%
\textsuperscript{\@fnsymbol{#1}}%
}
\makeatother


% ##############################################################################
% Numbering, so that the references read as in the IEEE style of `paper.md`:
% sections I, II, ..., subsections VII-B, tables I, II, ..., "Fig." in captions,
% and display equations numbered (1), (2), ... as in its PDF, so that they can be
% cited in a review.
% ##############################################################################

\renewcommand{\thesection}{\Roman{section}}
\renewcommand{\thesubsection}{\thesection-\Alph{subsection}}
\renewcommand{\thetable}{\Roman{table}}
\renewcommand{\figurename}{Fig.}

% Subsection titles on their own line, instead of the run-in titles of `amsart`.
\makeatletter
\def\subsection{\@startsection{subsection}{2}{\z@}%
  {.7\linespacing\@plus\linespacing}{.3\linespacing}%
  {\normalfont\bfseries}}
\makeatother

% "Proposition 1", numbered on its own and not within the section.
\theoremstyle{definition}
\newtheorem{prop}{Proposition}

% Numeric citations in order of first use, as in `paper.md`.
\usepackage[numbers,square]{natbib}

% The figures are found from the root of the repo, where `run_latex.py` runs
% `pdflatex`, and from the paper dir.
\graphicspath{{figures/}{papers/Optimal_strategy_for_racket_sports/figures/}}

% ##############################################################################
% Notes for GP.
% ##############################################################################

% Each open question sits where it concerns the text, in a colored box. Set
% `\gpnotesfalse` to hide all the notes, e.g., for the submission version. The
% notes have no `\cite`, so hiding them does not renumber the references.
\usepackage[most]{tcolorbox}
\newif\ifgpnotes
\gpnotestrue
\newcommand{\gpnote}[2]{%
  \ifgpnotes
    \begin{tcolorbox}[
      breakable,
      colback=yellow!8,
      colframe=orange!75!black,
      fonttitle=\bfseries\small,
      fontupper=\small,
      title={Note for GP: #1}]
      #2
    \end{tcolorbox}%
  \fi}

\begin{document}
  \title{Optimal Shot Placement in Racket Sports: A Game-Theoretic Framework}

  \author{Aahan Rembersu}
  \address{Department of Computer Science, University of Maryland, College
    Park, College Park, MD, USA}
  \email{arembers@terpmail.umd.edu}

  \author{Giacinto Paolo (GP) Saggese}
  \address{Department of Computer Science, University of Maryland, College
    Park, College Park, MD, USA}
  \email{gsaggese@umd.edu}

  \keywords{tennis, pickleball, shot placement, game theory, reinforcement
    learning, contextual bandit, Monte Carlo methods}

  \date{\today}

  \begin{abstract}
    Where to place a shot is a central tactical decision in tennis and
    pickleball. We model the return of serve in two dimensions and estimate by
    Monte Carlo, on a grid of aim points, the probability that the return lands
    in and the server cannot reach it. Four findings stand out.
    \begin{itemize}
      \item[(i)] \emph{The best return is central:} at the default speeds, it
      is the middle of the server's half in both sports, not the lines or the
      corners.

      \item[(ii)] \emph{Only the speed ratio matters:} with no reaction time and
      a spread independent of the ball speed, the winner probability depends on
      the two speeds only through their ratio.

      \item[(iii)] \emph{Wins and losses are enough to learn it:} a proximal
      policy optimization (PPO) agent loses under 0.8 percentage points of
      winner probability in expectation, against 11 to 21 for a random aim
      point, and its greedy aim falls in the best region on 43 of 54 serves.

      \item[(iv)] \emph{Hitting much faster does not pay when it costs
      accuracy:} if doubling the ball speed doubles the spread in depth ($a =
      1$, the strongest coupling we assume), a ball twice the default speed
      loses 9.9 points of the highest winner probability in pickleball and 11.1
      points of its mean in tennis.
    \end{itemize}
    We also describe the general framework: a trajectory model, a grid score
    with reachability, and a zero-sum game over placements, formulated but not
    solved.
  \end{abstract}

  \maketitle

  \gpnote{title (Gdoc Q19)}{%
    Keep ``A Game-Theoretic Framework'' in the title? The paper computes the 2D
    model, its Monte Carlo maps, and the reinforcement learning check; the game
    of Section~\ref{sec:game} is formulated but not solved. \emph{Our
    recommendation:} a title that names what is computed, e.g., ``Where to Aim
    the Return in Tennis and Pickleball: Monte Carlo and Reinforcement
    Learning'', with the game as future work. Your call.}

  \gpnote{abstract, finding (iv)}{%
    Finding (iv) is stated at $a = 1$ only, the end of our sweep. \#984
    computes the threshold $a^\star$ of the coefficient above which hitting
    faster than $v_0$ stops paying, per sport and shape, as you asked on Slack
    on 2026-11-17. Once it is merged, we will add $a^\star$ to
    Section~\ref{sec:speed} and turn (iv) into a threshold claim.}

  % ############################################################################
  \section{Introduction}\label{sec:intro}

  Deciding where to send the ball is, together with deciding how to hit it, the
  fundamental tactical choice in racket sports. A shot that lands safely in the
  court but travels straight to the opponent achieves little; a shot aimed at
  the extreme corner of the court is more likely to be a winner if it lands in,
  but is also more likely to land out.

  Optimal shot placement is therefore a decision under uncertainty, made under
  tight time constraints, and made against an opponent who reacts strategically
  to the pattern of one's choices.

  Prior work on this problem tends to fall into two categories. The first
  category relies on high-fidelity physical simulation of ball flight, typically
  in three dimensions and including spin and aerodynamic drag, combined with
  Monte Carlo sampling over the space of possible strokes \cite{cross2014drag}.
  Such simulations are accurate but computationally expensive, and their
  complexity makes it difficult to reason analytically about why one placement
  is preferable to another.

  The second category is data-driven: it uses ball- and player-tracking data to
  predict where a shot is likely to go \cite{wei2013shotlocations} or to
  evaluate the outcomes of shots that were actually played
  \cite{kovalchik2020voncramm}. These methods are descriptive rather than
  prescriptive; they characterize observed behavior but do not, by themselves,
  prescribe an optimal placement policy from first principles.

  This paper proposes a middle ground: a reduced-order, first-principles
  framework that is simple enough to reason about analytically, yet rich enough
  to capture the trade-off between shot accuracy and opponent reachability that
  we regard as the crux of shot-placement strategy.

  We compute with its simplest instance, a two-dimensional model of the return
  of serve, which rests on three simplifications.

  \begin{itemize}
    \item The flight of the ball is not simulated. The return lands at a random
    point, drawn from a Gaussian around the aim point, and flies to it in a
    straight line at a constant speed. The trajectory model of
    Section~\ref{sec:trajectory} is one way to derive this spread from the
    errors of the stroke.

    \item The court is discretized into a grid of aim points, and the
    probability that a return aimed at each one is a winner is estimated by
    Monte Carlo, with the same random draws for every aim point.

    \item The server reaches the ball if, after a fixed reaction time, they can
    run to the landing point at a constant top speed before the ball lands. This
    is checked for each sampled landing point, not as a separate probability.
  \end{itemize}

  The main result runs against a natural intuition, which is to aim a return
  near the lines and corners, far from the opponent.
  \gpnote{corners claim (Gdoc Q20)}{%
    Keep ``a natural intuition is to aim near the lines and corners'' as an
    intuition, or add a source? The main result is contrasted with it, and the
    bibliography has no source for where players aim their returns; your
    outreach hook (Gdoc Q16) is ``pros go for the corners''. \emph{Our
    recommendation:} keep it as an intuition, unless you have a source in
    mind.}
  At the default speeds of the model, whose reaction time and player speed are
  within the ranges reported in the literature
  (Section~\ref{sec:assumptions}), the best return is instead the middle of the
  server's half, in both sports (Fig.~\ref{fig:winner-map}): a return near a
  line is more likely to land out, and the server covers too little ground
  before the ball lands for distance from them to pay for that risk. A
  reinforcement learning agent that sees only whether each return was a winner
  learns the same answer (Section~\ref{sec:rl}).

  The main contributions of this paper are:

  \begin{itemize}
    \item A two-dimensional model of the return of serve, with a closed form for
    the in-bounds probability and a Monte Carlo estimate of the winner
    probability as one joint event, on common random numbers, with the best
    placement reported as a region of statistically tied aim points
    (Sections~\ref{sec:pin} and~\ref{sec:pwin}).

    \item A study of which variables change the answer, with a proposition
    that, with no reaction time, the ball and player speeds act only through
    their ratio (Section~\ref{sec:variables} and Appendix~A), and of the best
    ball speed when a faster return is less accurate
    (Section~\ref{sec:speed}).

    \item A one-step reinforcement learning environment built on the same
    model, which reproduces the winner probability map from the reward alone,
    and a proximal policy optimization (PPO) agent evaluated on it with a pass
    rule written before the run (Section~\ref{sec:rl}).

    \item The general framework that the 2D model instantiates: a 1D/2D
    reduction of the trajectory-feasibility problem with an error-propagation
    model (Section~\ref{sec:problem}), a grid score that combines in-bounds
    probability with a deterministic reachability indicator
    (Sections~\ref{sec:discretization} to~\ref{sec:complexity}), and a zero-sum
    game over placements connected to empirical evidence of minimax play at the
    professional level \cite{walker2001minimax} (Section~\ref{sec:game}).

    \item A demonstration that tennis and pickleball can be described by the
    same parametrized model, differing primarily in maximum ball speed and court
    scale, illustrated with a small worked example that exposes qualitatively
    different optimal regimes for the two sports (Sections~\ref{sec:example}
    and~\ref{sec:sports}).
  \end{itemize}

  The paper is organized as follows. Section~\ref{sec:related} reviews related
  work. Section~\ref{sec:problem} formalizes the problem.
  Section~\ref{sec:grid} describes the grid-based scoring methodology and, in
  Sections~\ref{sec:pin} and~\ref{sec:pwin}, the 2D model we compute with.
  Section~\ref{sec:game} develops the game-theoretic layer.
  Section~\ref{sec:example} works a small illustrative example by hand.
  Section~\ref{sec:sports} discusses the tennis and pickleball
  parametrizations, which variables matter, and the best ball speed when speed
  costs accuracy. Section~\ref{sec:rl} checks the return of serve with
  reinforcement learning. Section~\ref{sec:discussion} discusses limitations,
  and Section~\ref{sec:conclusion} concludes. Appendix~B lists, for each
  result, the notebook that reproduces it and the entry of the project's
  research log that holds its runs and figures.

  % ############################################################################
  \section{Related Work}\label{sec:related}

  \textbf{Ball trajectory physics.} The aerodynamics of racket-sport balls,
  including drag and Magnus-effect lift due to spin, have been measured directly
  using high-speed cameras and wind-tunnel experiments \cite{cross2014drag}.
  These measurements show that drag and lift forces are not negligible relative
  to gravity, particularly at the higher speeds reached by serves and hard
  groundstrokes. Our framework deliberately ignores these effects to keep the
  trajectory problem analytically tractable, at the cost of introducing a
  systematic bias that we return to in Section~\ref{sec:discussion}.

  \textbf{Data-driven shot analytics.} Using Hawk-Eye-style tracking data, prior
  work has built probabilistic models that predict the location of a player's
  next shot from the spatiotemporal features of the current rally
  \cite{wei2013shotlocations}, and has modeled the space of plausible future
  shot trajectories to evaluate the quality of a player's decision relative to
  the alternatives available at the time \cite{kovalchik2020voncramm}. These
  methods answer ``what shot is this player likely to play, and how good was
  it,'' which is a different question from ``what shot should any player play
  given the current geometry,'' the question addressed here. Our framework does
  not require historical tracking data and can, in principle, be evaluated for
  any court geometry and any pair of player positions.

  \textbf{Game theory in racket sports.} The clearest empirical link between
  game theory and tennis strategy is the finding that the serve-direction
  choices of elite players are statistically consistent with mixed-strategy
  minimax equilibrium: players who lose a disproportionate share of points on
  serves to one side subsequently increase the frequency of serves to the other
  side, and the win probabilities across serve directions are close to equal,
  as predicted by minimax play \cite{walker2001minimax}. That result concerns a
  small, discrete choice, direction of serve, left, center, or wide.
  Section~\ref{sec:game} extends the same equilibrium logic to a continuous grid
  of target cells and to shots beyond the serve.

  \textbf{Reinforcement learning.} Reinforcement learning has been applied to
  racket sports before. Mohan~\cite{mohan2025tennis} trains a dueling double
  deep Q-network on a custom tennis simulator to choose among ten rally-level
  tactical actions, and Ding et al.~\cite{ding2022racket} learn a Q-function
  from player poses and shuttlecock positions to evaluate badminton players in
  world-class matches. In both, the learned values have no independent answer
  to be checked against. Here the in-bounds probability has a closed form,
  checked by Monte Carlo, and the learned policy is validated against a brute
  force over every aim point (Section~\ref{sec:rl}): reinforcement learning is
  a check of the model, not the source of the answer. A single return of serve
  is a contextual bandit, a one-step decision with a context, the serve, and a
  reward \cite{lattimore2020bandit}. We state it as a Gymnasium environment
  \cite{towers2024gymnasium} and train PPO \cite{schulman2017ppo}, as
  implemented in Stable-Baselines3 \cite{raffin2021sb3}. Extending the
  environment to rallies and to an opponent who adapts leads back to the game
  of Section~\ref{sec:game}.

  % ############################################################################
  \section{Problem Formulation}\label{sec:problem}

  \subsection{Court Geometry and Notation}\label{sec:court}

  We model the court as a planar rectangle and place the origin at the
  striker's contact point.

  Table~\ref{tab:court} summarizes the court dimensions and net heights defined
  by the respective governing bodies, together with the fastest recorded shot
  speed in each sport, which we use as an upper bound on the maximum ball speed
  parameter $v_{\max}$.

  \begin{table}[htbp]
    \centering
    \caption{Court and ball-speed parameters for tennis and pickleball}
    \label{tab:court}
    \begin{tabular}{lrr}
      \toprule
      Parameter & Tennis & Pickleball \\
      \midrule
      Court (singles), length x width & 23.77 m x 8.23 m & 13.41 m x 6.10 m \\
      Net height, center / post & 0.914 m / 1.07 m & 0.86 m / 0.91 m \\
      Non-volley zone (from net) & n/a & 2.13 m \\
      Fastest recorded shot & 263 km/h (73.1~m/s) & 99.8 km/h (27.7~m/s) \\
      \bottomrule
    \end{tabular}
  \end{table}

  Tennis dimensions and net heights are from the ITF Rules of Tennis
  \cite{itf2024rules}; pickleball dimensions are from the USA Pickleball
  official rulebook \cite{usap2025rulebook}; the fastest-shot figures are the
  certified records for a men's serve \cite{guinness2012serve} and for a
  hard-hit shot in pickleball \cite{guinness2025pickleball}, respectively.

  \subsection{Simplifying Assumptions}\label{sec:assumptions}

  The framework rests on the following approximations, which trade physical
  fidelity for tractability:

  \begin{itemize}
    \item \textbf{Point-mass ball and point players.} Racket or paddle contact
    is modeled as an instantaneous velocity change at a point, and players are
    modeled as points on the court plane rather than as extended bodies with
    reach.

    \item \textbf{No air drag or spin.} Ball flight follows unpowered projectile
    motion. This is the most consequential simplification;
    Section~\ref{sec:discussion} discusses its likely impact using measured
    drag and lift coefficients \cite{cross2014drag}.

    \item \textbf{Bounded reaction time.} A player requires a fixed latency
    after ball contact before beginning to move. We use the 190 to 250 ms range
    reported for simple visual reaction time as a plausible baseline
    \cite{woods2015reaction}, noting that this can be a conservative estimate:
    anticipation and cue reading can shorten the effective decision latency in
    practice, while decision complexity can lengthen it.

    \item \textbf{Bounded movement speed.} A player's on-court movement speed is
    bounded by a sport- and player-specific maximum $v_p$. Measured average
    movement speeds immediately before and after the split-step range from
    about 0.55 to 1.9~m/s depending on stroke type and player level
    \cite{filipcic2017splitstep}; peak short-burst court-coverage speed is
    higher, and we treat $v_p$ as a free parameter of the model rather than
    fixing it to a single literature value.

    \item \textbf{Groundstrokes only.} We model rallies in which each shot
    bounces before the opponent responds, and exclude volleys, in which a player
    strikes the ball before it bounces. Extending the framework to volleys is
    left to future work.
  \end{itemize}

  \subsection{Problem 1: Trajectory Feasibility}\label{sec:trajectory}

  For a given shot, let:

  \begin{itemize}
    \item $S$ denote the striker's contact point
    \item $N$ the point where the trajectory crosses the net, and
    \item $T$ a candidate target on the opponent's side of the court.
  \end{itemize}

  Because $S$, $N$, and $T$ are, by construction, coplanar for a fixed aim
  direction, we work in the vertical plane containing the line $ST$, with
  horizontal coordinate $x$ (distance along $ST$) and vertical coordinate $z$
  (height above the court surface). Lateral placement, i.e., the choice of aim
  direction itself, is set by the choice of target cell on the court grid
  (Section~\ref{sec:discretization}); its execution error is the lateral angle
  $\phi$ introduced below.

  Given a contact point $S$ at height $h_0$, a net at horizontal distance
  $d_{\mathrm{net}}$ and height $h_{\mathrm{net}}$, and a target $T$ at
  horizontal distance $d_T$, the trajectory-feasibility problem asks for a
  launch speed $v_0$ and launch angle $\theta$ (measured from horizontal) such
  that the ball clears the net and lands at $T$. Under the no-drag assumption,
  the trajectory in the vertical plane is
  \begin{equation}
  x(t) = v_0 \cos\theta \, t, \qquad
  z(t) = h_0 + v_0 \sin\theta \, t - \frac{1}{2} g t^2 .
  \end{equation}
  The flight time to landing ($z = 0$) is
  \begin{equation}
  T_f = \frac{v_0 \sin\theta + \sqrt{(v_0 \sin\theta)^2 + 2 g h_0}}{g} ,
  \end{equation}
  and the net-clearance constraint requires $z(t_{\mathrm{net}}) >
  h_{\mathrm{net}}$ at $t_{\mathrm{net}} = d_{\mathrm{net}} / (v_0
  \cos\theta)$. The feasible set $\mathcal{F}(T)$ for hitting $T$ within a
  tolerance $\epsilon$ then consists of every pair $(v_0, \theta)$ that
  satisfies both the net-clearance condition above and the landing condition
  \begin{equation}
  x(T_f) \in [\, d_T - \epsilon, \; d_T + \epsilon \,] ,
  \end{equation}
  subject to $v_0 \le v_{\max}$ for the sport in question
  (Table~\ref{tab:court}). Because this is a two-parameter feasibility problem
  posed along a single spatial axis, we refer to it as the ``1D'' trajectory
  problem, in contrast to a full 3D flight simulation.

  A player cannot execute $(v_0, \theta)$ exactly; every stroke carries
  execution noise. We model this noise as
  \begin{equation}
  \theta = \bar\theta + \delta_\theta, \qquad
  \delta_\theta \sim \mathcal{N}(0, \sigma_\theta^2),
  \end{equation}
  \begin{equation}
  v_0 = \bar v_0 (1 + \delta_v), \qquad
  \delta_v \sim \mathcal{N}(0, \sigma_v^2),
  \end{equation}
  where $\bar\theta$ and $\bar v_0$ are the intended (aim-point-feasible) angle
  and speed. A third error term, the lateral angle $\phi \sim \mathcal{N}(0,
  \sigma_\phi^2)$, perturbs the aim direction within the horizontal plane and
  is what makes the landing distribution genuinely two-dimensional rather than
  confined to the line $ST$. Propagating $(\delta_\theta, \delta_v, \phi)$
  through the trajectory equations, either by direct Monte Carlo sampling or by
  local linearization for small errors, yields a distribution over landing
  points on the court plane, which we use in Section~\ref{sec:grid} to score
  candidate targets.

  \subsection{Problem 2: Placement Selection}\label{sec:placement}

  Given the current ball state at the moment of contact at the attacking
  player's position $S$ and the returning opponent's court position $O$, the
  placement-selection problem asks for the target $T$ that maximizes the chance
  the shot is both a legal, in-bounds shot and one the opponent cannot reach
  before it bounces twice. This requires combining the landing distribution
  from Section~\ref{sec:trajectory} with a model of the opponent's response
  capability, which we formalize next.

  The computations of this paper use a two-dimensional reduction of both
  problems, defined in Sections~\ref{sec:pin} and~\ref{sec:pwin}. The landing
  point is drawn directly from a Gaussian around the aim point, instead of
  propagating stroke errors through the trajectory of Problem 1, and the
  opponent's reach is checked for each sampled landing point.
  Sections~\ref{sec:sports} and~\ref{sec:rl} study this model.

  % ############################################################################
  \section{Grid-Based Reachability and Monte Carlo Scoring}\label{sec:grid}

  \subsection{Court Discretization}\label{sec:discretization}

  We discretize the opponent's court into a grid of candidate target cells
  $\mathcal{C} = \{c_1, \dots, c_M\}$. Grid resolution trades accuracy for
  computational cost: a finer grid better approximates the continuous placement
  decision but requires proportionally more Monte Carlo samples and more
  reachability evaluations, discussed in Section~\ref{sec:complexity}.

  \subsection{Landing Probability via Monte Carlo}\label{sec:landing}

  For a candidate cell $c$ with center $T_c$, we compute the
  feasibility-constrained aim parameters $(\bar\theta, \bar v_0)$ from
  Section~\ref{sec:trajectory}, draw $K$ samples of $(\delta_\theta, \delta_v,
  \phi)$ from the error model, simulate the resulting trajectory for each
  sample, and estimate
  \begin{equation}
  \begin{aligned}
  P_{\mathrm{in}}(c) \approx \frac{1}{K} \sum_{k=1}^{K} \mathbb{1}[\,
  & z_k(t_{\mathrm{net}}) > h_{\mathrm{net}} \\
  & \wedge \mathrm{land}(\theta_k, v_{0,k}, \phi_k) \in \mathcal{L} \,] ,
  \end{aligned}
  \end{equation}
  where $\mathcal{L}$ is the legal landing region (the opponent's half in a
  rally, the service box for a serve) and $z_k(t_{\mathrm{net}}) >
  h_{\mathrm{net}}$ is the net-clearance condition of
  Section~\ref{sec:trajectory} for sample $k$. A sample counts if the ball
  clears the net and lands anywhere in $\mathcal{L}$: the cell $c$ is where the
  striker aims, not where the ball must land, so $P_{\mathrm{in}}(c)$ is the
  in-bounds probability of a shot aimed at $c$. The feasible set
  $\mathcal{F}(T_c)$ holds many launch pairs: we evaluate a grid of feasible
  launch angles $\bar\theta$ for each cell, each with the launch speed that
  lands it at $T_c$, and keep for each cell the angle with the highest score
  $S(c)$ of Section~\ref{sec:score}, since a higher, slower shot flies longer
  and the best angle depends on the opponent's position.

  Because the underlying trajectory model is one-dimensional per sample
  (Section~\ref{sec:trajectory}), each Monte Carlo draw costs a closed-form
  evaluation rather than a numerical integration of drag-augmented equations of
  motion, which is what makes $K$ in the thousands practical per cell even for
  moderately fine grids.

  \subsection{Deterministic Reachability}\label{sec:reach}

  Let $t_r$ denote the opponent's reaction time and $v_p$ the opponent's maximum
  movement speed (Section~\ref{sec:assumptions}). For a candidate cell $c$ at
  distance $d(O, c)$ from the opponent's current position $O$, we define
  reachability as the deterministic indicator
  \begin{equation}
  R(c) = \mathbb{1}\big[\, t_r + d(O,c)/v_p \; \le \; T_f(c) \,\big] ,
  \end{equation}
  where $T_f(c)$ is the ball's flight time to $c$ from
  Section~\ref{sec:trajectory}. Unlike $P_{\mathrm{in}}$, $R$ is not estimated
  by sampling: given fixed $t_r$, $v_p$, and $T_f(c)$, whether the opponent's
  minimum possible arrival time is within the flight-time budget is a yes/no
  geometric fact, not a random outcome. Randomness in the opponent's true
  reaction time or top speed can be incorporated by treating $R(c)$ as a
  probability rather than an indicator; we retain the deterministic form here
  for tractability and note the extension in Section~\ref{sec:discussion}.

  \subsection{Composite Score and Candidate Selection}\label{sec:score}

  We combine the two quantities into a single score per cell,
  \begin{equation}
  S(c) = P_{\mathrm{in}}(c) \cdot \big(1 - R(c)\big) ,
  \end{equation}
  which is zero for any cell the opponent can reach, regardless of how safe the
  shot is, and equal to the in-bounds probability for any cell the opponent
  cannot reach. The candidate placement is then
  \begin{equation}
  c^\star = \argmax_{c \in \mathcal{C}} S(c) .
  \end{equation}
  This multiplicative form directly encodes the qualitative intuition that a
  placement is only as good as the product of ``does it land in'' and ``can the
  opponent not get there.'' Section~\ref{sec:example} works through a concrete
  instance of this scoring rule. Weighted or non-multiplicative combinations of
  $P_{\mathrm{in}}$ and $R$ are possible, for example to express risk
  aversion, but are left to future work.

  \subsection{Serve and Return}\label{sec:serve}

  The same grid-and-score procedure applies to serving and to returning serve,
  with two changes to the setup: the target grid is restricted to the legal
  service box, and $v_{\max}$ and the tolerance $\epsilon$ reflect the serve
  motion rather than a groundstroke. The opponent's position $O$ for a return
  is the returner's ready position rather than a mid-rally recovery position,
  and $t_r$, $v_p$ can differ from rally-play values because return positioning
  and split-step timing are typically more standardized than in-rally movement.

  \subsection{Computational Complexity}\label{sec:complexity}

  Evaluating one grid consists of $M$ cells, each requiring $K$ trajectory
  samples for $P_{\mathrm{in}}$ and one closed-form reachability check for $R$,
  giving $O(MK)$ closed-form trajectory evaluations per decision. This is
  substantially cheaper than a full 3D drag-and-spin simulation integrated
  numerically over the same $M \times K$ grid, because each sample here is an
  $O(1)$ evaluation of the closed-form no-drag equations of
  Section~\ref{sec:trajectory} rather than a numerical solution of an ordinary
  differential equation. The reduction from 3D to a per-cell 1D/2D trajectory
  model is therefore what makes exhaustive grid scoring, rather than a small
  number of hand-picked candidate shots, computationally practical within the
  time constraints of a rally.

  \subsection{Closed-Form In-Bounds Probability for a Gaussian Landing
    Point}\label{sec:pin}

  Section~\ref{sec:landing} estimates $P_{\mathrm{in}}$ by propagating stroke
  errors through the trajectory model. The model we compute with is simpler:
  it places the uncertainty directly on the landing point. A shot aimed at
  $(x_2, y_2)$ lands at $(X, Y)$ with
  \begin{equation}
  X \sim \mathcal{N}(x_2, \sigma_x^2), \qquad
  Y \sim \mathcal{N}(y_2, \sigma_y^2),
  \end{equation}
  where $X$ and $Y$ are independent. This model uses a frame with the origin at
  the center of the net, $x$ across the court and $y$ along it, rather than the
  striker-centered frame of Section~\ref{sec:court}. For a court of length $L$
  and width $W$, a shot toward the half $y \le 0$ is in when $|X| \le W/2$ and
  $-L/2 \le Y \le 0$. Because the two coordinates are independent and this
  region is a rectangle, the event splits into one condition on $X$ and one on
  $Y$, and $P_{\mathrm{in}}$ is the product of two differences of the standard
  normal cumulative distribution function (CDF) $\Phi$:
  \begin{equation}
  \begin{aligned}
  P_{\mathrm{in}} &= p_x \, p_y , \\
  p_x &= \Phi\left(\frac{W/2 - x_2}{\sigma_x}\right)
  - \Phi\left(\frac{-W/2 - x_2}{\sigma_x}\right) , \\
  p_y &= \Phi\left(\frac{-y_2}{\sigma_y}\right)
  - \Phi\left(\frac{-L/2 - y_2}{\sigma_y}\right) ,
  \end{aligned}
  \end{equation}
  where $p_x$ is the probability that $X$ lands between the sidelines and $p_y$
  the probability that $Y$ lands between the baseline and the net.

  With correlated errors this would require a bivariate normal CDF, and with a
  non-rectangular region a two-dimensional integral. The model ignores the net:
  a shot counts as out only if it lands outside the target half.

  In the accompanying notebook, the serve and return setting (compare
  Section~\ref{sec:serve}) is modeled as follows. The server chooses the deuce
  or the ad side and stands behind the baseline on that side's half of the
  court. Only the service box diagonally across the court is legal, and lines
  are in, except the pickleball non-volley-zone line. A serve that is wide,
  lands in the wrong box, lands long, or, in pickleball, lands in the
  non-volley zone is a fault, and a fault ends the point with no return. After
  a legal serve to $(x_1, y_1)$, the returner moves to $(x_1, y_1)$ and returns
  toward $(x_2, y_2)$ with spread $(\sigma_x, \sigma_y)$, and $P_{\mathrm{in}}$
  is the probability that the return lands in the server's half. The closed
  form is checked against a Monte Carlo estimate from samples of the same
  Gaussian, which agrees with it within sampling error. It is implemented, with
  an interactive notebook for both sports, in the accompanying code
  (Appendix~B).

  Details: research log, Experiment 1 (Appendix~B).

  \subsection{Winner Probability of a Return}\label{sec:pwin}

  This subsection extends Section~\ref{sec:pin} from ``does the return land
  in'' to ``is the return a winner'': it lands in and the server cannot reach
  it. After a legal serve to $(x_1, y_1)$, the returner hits from $(x_1, y_1)$
  toward $(x_2, y_2)$ with the Gaussian spread of Section~\ref{sec:pin}, and
  the server waits at the center of their baseline, $P = (0, -L/2)$. For each
  of $K$ sampled landing points $(X_k, Y_k)$, the flight time is the
  straight-line distance over a constant ground speed $v_b$ of the ball,
  \begin{equation}
  T_k = \frac{\lVert (X_k, Y_k) - (x_1, y_1) \rVert}{v_b} ,
  \end{equation}
  so it depends on where the ball lands, not on where it was aimed. The server
  reacts for $t_r$ and then runs in a straight line at speed $v_p$, and reaches
  the landing point if
  \begin{equation}
  \lVert (X_k, Y_k) - P \rVert \le v_p \max(0, \, T_k - t_r) .
  \end{equation}
  Unlike Section~\ref{sec:reach}, the reach is evaluated for each sample. Let
  $W_k$ be the event that sample $k$ lands in the server's half $H$ and is not
  reached. The winner probability is estimated as this one joint event on the
  same samples,
  \begin{equation}
  P_{\mathrm{win}}(x_2, y_2) \approx \frac{1}{K} \sum_{k=1}^{K}
  \mathbb{1}[W_k] .
  \end{equation}
  It is not computed as the product $P_{\mathrm{in}} \cdot P(\text{not
  reached})$, since both events depend on the same landing point. The model is
  two-dimensional: the ball flies in a straight line at a constant speed, the
  bounce is ignored, and the server must reach the landing point by the time
  the ball lands.

  $P_{\mathrm{win}}$ is evaluated at the centers of a grid over the server's
  half ($10 \times 15$ cells for tennis and $9 \times 10$ for pickleball, with
  $K = 5000$), with common random numbers: every target uses the same standard
  normal draws, shifted to its own aim point, so two targets differ because of
  where they aim rather than because of sampling noise. Many targets can tie at
  the top, for example when all $K$ samples are winners, so the best placement
  is reported as a region: every target within two standard errors of the
  highest $P_{\mathrm{win}}$, using the standard error of the difference of two
  estimates, marked at the cell closest to the region's centroid. Each map is
  identified by a settings string that lists all its inputs: the sport, the
  side, the serve, the spread, $v_b$, $t_r$, $v_p$, the seed, the grid size,
  and $K$, so a map reported in the research log can be rebuilt exactly in the
  notebook.

  With a ball speed of 25~m/s for tennis, $t_r = 0.2$ s, and $v_p = 1.5$~m/s, a
  return that lands next to the server flies about 15 m in 0.61 s, so the
  server covers only about 0.6 m. $P_{\mathrm{win}}$ is then close to
  $P_{\mathrm{in}}$ except in a small dip around the server, and the best
  region is the middle of the server's half rather than a spot near the lines
  (Fig.~\ref{fig:winner-map}). The region is large, 41 of 150 cells, because
  these cells tie within sampling error near $P_{\mathrm{win}} = 1$: placement
  inside it hardly matters. Because the spread does not depend on the ball
  speed in this model, a faster return never lowers $P_{\mathrm{win}}$, so
  hitting as fast as possible always looks best; Section~\ref{sec:speed} makes
  the spread grow with the ball speed. The model is implemented, with a heatmap
  of $P_{\mathrm{win}}$ for both sports, in the accompanying code
  (Appendix~B).

  \begin{figure}[htbp]
    \centering
    \includegraphics[width=0.42\textwidth]{winner_heatmap_tennis.png}
    \caption{$P_{\mathrm{win}}$ of every aim point for tennis at the default
      settings: serve at $(-2.1, 3.2)$, spread $(0.5, 1.0)$ m, $v_b = 25$~m/s,
      $t_r = 0.2$ s, $v_p = 1.5$~m/s, and $K = 5000$. The black outline is the
      best region, within two standard errors of the highest value (41 of 150
      cells), and the star is its middle: the best return is the middle of the
      server's half, not the lines or the corners.}
    \label{fig:winner-map}
  \end{figure}

  Details: research log, Experiment 2 (Appendix~B).

  % ############################################################################
  \section{Game-Theoretic Placement}\label{sec:game}

  \subsection{Placement as a Zero-Sum Game}\label{sec:zerosum}

  Section~\ref{sec:grid} treats the opponent's position $O$ as fixed and known,
  and selects the cell that maximizes $S(c)$ against that fixed position. A
  rational opponent, however, anticipates the striker's tendencies and adjusts
  $O$, the pre-shot recovery position, accordingly. If the striker always aimed
  at the single highest-scoring cell, an opponent who learns this would simply
  recover toward that cell, driving its effective $R(c)$ toward 1 and its score
  toward zero. Optimal play therefore requires mixing across multiple candidate
  cells rather than deterministically selecting the argmax.

  We formalize repeated placement decisions as a zero-sum game in which the
  striker chooses a mixed strategy $\sigma \in \Delta(\mathcal{C})$ over target
  cells and the returner chooses a mixed strategy $\rho \in \Delta(\mathcal{C})$
  over anticipated recovery positions, with payoff $u(c, c')$ to the striker
  equal to the point-win probability when the shot is aimed at $c$ and the
  returner has recovered toward $c'$. The equilibrium strategy is
  \begin{equation}
  \sigma^\star = \argmax_{\sigma \in \Delta(\mathcal{C})} \;
  \min_{\rho \in \Delta(\mathcal{C})} \;
  \EE_{c \sim \sigma,\, c' \sim \rho}[\, u(c, c') \,] .
  \end{equation}
  For a grid with $M$ cells, this reduces to solving an $M \times M$ zero-sum
  matrix game, which is a linear program and is computationally far cheaper
  than the $O(MK)$ scoring step that produces the payoff matrix in the first
  place.

  We do not solve this game in this paper. The computations of
  Sections~\ref{sec:sports} and~\ref{sec:rl} keep the server at a fixed
  position; extending them to a server who adapts, for example by self-play, is
  future work (Section~\ref{sec:conclusion}).

  \subsection{Connection to Minimax Play at Wimbledon}\label{sec:wimbledon}

  This formulation generalizes the discrete, three-direction (left, center,
  wide) serve-placement game empirically validated by Walker and Wooders
  \cite{walker2001minimax}, who found that professional servers' direction
  frequencies and the resulting win probabilities across directions are
  statistically indistinguishable from minimax equilibrium play. Their result
  concerns a fixed, small strategy space chosen by convention; the framework
  proposed here suggests that the same equilibrium logic should, in principle,
  extend to the full continuous target grid and to shots beyond the serve,
  provided the payoff matrix $u(c, c')$ can be estimated. Testing this
  generalization against real rally data is a natural direction for future work
  (Section~\ref{sec:conclusion}).

  \subsection{Post-Shot Positioning}\label{sec:postshot}

  A related and still-open question is where the striker should move
  immediately after hitting the shot, before knowing the opponent's response.
  This is the mirror image of the placement problem: having committed to a
  distribution $\sigma$ over target cells, the striker's own recovery position
  should maximize coverage against the returner's best-response distribution
  over the striker's own court. In general this defines a bilevel, repeated
  game in which each player's optimal positioning depends on the other's
  strategy at the following shot. It can, in principle, be approximated with
  the same grid representation and standard zero-sum solution methods (e.g.,
  the double oracle algorithm), applied recursively shot by shot.

  % ############################################################################
  \section{Illustrative Worked Example}\label{sec:example}

  To make the scoring rule of Section~\ref{sec:grid} concrete, this section
  works through a small, hand-computed example. It illustrates the scoring rule
  by hand and is not a simulated experiment; the computed results use the 2D
  model of Sections~\ref{sec:pin} and~\ref{sec:pwin} and are in
  Sections~\ref{sec:sports} and~\ref{sec:rl}.

  \subsection{Setup}\label{sec:setup}

  Fig.~\ref{fig:trajectory} shows the trajectory-feasibility geometry
  underlying the example: a contact point $S$ at height $h_0 = 1.0$ m, a net at
  $d_{\mathrm{net}} = 12$ m with $h_{\mathrm{net}} = 0.914$ m, and a target $T$
  at $d_T = 20$ m. For a launch angle $\theta = 8^\circ$ and speed $v_0 =
  22.9$~m/s, the resulting trajectory clears the net by about $0.40$ m and lands
  at $T$ after a flight time $T_f = 0.88$ s; the dashed curves show four
  perturbed trajectories under $\pm 1.5^\circ$ and $\pm 3^\circ$ angle error
  combined with $\mp 5$ to $8\%$ speed error, and the shaded segment shows the
  resulting spread of landing points.

  \begin{figure}[htbp]
    \centering
    \includegraphics[width=0.62\textwidth]{trajectory_schematic.png}
    \caption{1D trajectory-feasibility geometry: contact point $S$, net, target
      $T$, the nominal trajectory (solid), four perturbed trajectories
      (dashed), and the resulting landing spread.}
    \label{fig:trajectory}
  \end{figure}

  We consider five candidate target cells $c_1, \dots, c_5$ at increasing
  lateral distance $d(O, c_i) \in \{0.3,\allowbreak 0.7,\allowbreak 1.1,\allowbreak 1.5,\allowbreak 1.9\}$ m from the
  returner's recovery position $O$, with illustrative in-bounds probabilities
  $P_{\mathrm{in}}(c_i) \in \{0.95,\allowbreak 0.91,\allowbreak 0.85,\allowbreak 0.77,\allowbreak 0.65\}$ decreasing with
  distance from center, consistent with the intuition that more extreme targets
  carry greater execution risk. We use round, illustrative values for the
  reachability parameters, chosen to be consistent with the ranges discussed in
  Section~\ref{sec:assumptions} rather than fitted to data: reaction time $t_r
  = 0.2$ s and movement speed $v_p = 1.5$~m/s.

  \subsection{Two Regimes}\label{sec:regimes}

  \textbf{Tennis groundstroke exchange.} With $T_f = 0.88$ s, the flight time of
  Fig.~\ref{fig:trajectory}, the available movement time after reacting is $T_f
  - t_r = 0.68$ s, giving a reach radius $v_p (T_f - t_r) = 1.02$ m. Cells $c_1$
  and $c_2$ fall within this radius ($R = 1$); cells $c_3$ to $c_5$ do not ($R =
  0$). Table~\ref{tab:toy} reports $P_{\mathrm{in}}$, $R$, and the composite
  score $S$ for each cell; Fig.~\ref{fig:toy} (left panel) shows the same values
  graphically. The maximizing cell is $c_3$, with $S(c_3) = 0.85$: not the
  safest cell ($c_1$, which is reachable and therefore scores zero) and not the
  most extreme cell ($c_5$, which is unreachable but carries the lowest
  in-bounds probability), but the closest-to-center cell that is just outside
  the opponent's reach envelope.

  \textbf{Pickleball net exchange.} Pickleball rallies at the non-volley-zone
  line involve much shorter ball-travel distances than tennis groundstrokes. As
  an illustration, a fast exchange over a $5$ m distance at $25$~m/s, a firm but
  sub-record pace relative to the certified $27.7$~m/s maximum in
  Table~\ref{tab:court}, gives $T_f = 0.2$ s. Because $T_f \le t_r$, the
  available movement time is zero for every cell, so $R(c_i) = 0$ for all $i$
  and $S(c) = P_{\mathrm{in}}(c)$ everywhere (Table~\ref{tab:toy},
  Fig.~\ref{fig:toy}, right panel). The maximizing cell is simply the safest
  one, $c_1$. This qualitative difference, reachability-driven trade-offs in
  tennis versus accuracy-dominated selection in fast pickleball exchanges,
  arises entirely from the ball-speed and distance parameters of
  Table~\ref{tab:court}, with the scoring mechanism of Section~\ref{sec:grid}
  held fixed.

  \begin{table}[htbp]
    \centering
    \caption{Toy-example scores for the five candidate cells under the tennis
      and pickleball parametrizations. Bold marks the argmax cell in each
      regime.}
    \label{tab:toy}
    \begin{tabular}{lrrrrrr}
      \toprule
      Cell & $d(O,c)$ (m) & $P_{\mathrm{in}}$ & $R$ (tennis) & $S$ (tennis)
        & $R$ (pball) & $S$ (pball) \\
      \midrule
      $c_1$ & 0.3 & 0.95 & 1 & 0.00 & 0 & \textbf{0.95} \\
      $c_2$ & 0.7 & 0.91 & 1 & 0.00 & 0 & 0.91 \\
      $c_3$ & 1.1 & 0.85 & 0 & \textbf{0.85} & 0 & 0.85 \\
      $c_4$ & 1.5 & 0.77 & 0 & 0.77 & 0 & 0.77 \\
      $c_5$ & 1.9 & 0.65 & 0 & 0.65 & 0 & 0.65 \\
      \bottomrule
    \end{tabular}
  \end{table}

  \begin{figure}[htbp]
    \centering
    \includegraphics[width=0.95\textwidth]{toy_example_scores.png}
    \caption{Toy-example scores for five candidate target cells, tennis (left)
      versus a fast pickleball net exchange (right). Gray bars: in-bounds
      probability $P_{\mathrm{in}}$; colored bars: composite score $S$; star:
      argmax cell.}
    \label{fig:toy}
  \end{figure}

  % ############################################################################
  \section{Tennis and Pickleball: A Unified Parametrization}\label{sec:sports}

  The equations in Sections~\ref{sec:problem} to~\ref{sec:game} make no
  sport-specific assumption beyond the numeric values of $v_{\max}$, court
  dimensions, and net height collected in Table~\ref{tab:court}. Formally, both
  sports are instances of the same model $\mathcal{M}(v_{\max},
  \text{geometry}, \sigma_\theta, \sigma_v, \sigma_\phi, t_r, v_p)$, and
  Section~\ref{sec:example} shows that the same scoring mechanism, applied with
  pickleball-consistent parameters instead of tennis-consistent ones, produces
  a qualitatively different optimal regime without any change to the model
  structure. The practical differences between the two sports, a smaller court,
  a lower net, a shorter or absent non-volley approach zone, and a paddle
  rather than a strung racket, enter the model exclusively through $v_{\max}$
  and the geometric constants of Table~\ref{tab:court}, rather than through
  additional equations. This supports treating ``sport'' as a parameter setting
  rather than a structural choice, as the computations of
  Sections~\ref{sec:variables} and~\ref{sec:rl} do. In the 2D model of
  Sections~\ref{sec:pin} and~\ref{sec:pwin}, the landing spread $(\sigma_x,
  \sigma_y)$ takes the place of the stroke errors $(\sigma_\theta, \sigma_v,
  \sigma_\phi)$, and the ground speed $v_b \le v_{\max}$ of the return takes
  the place of the launch speed, so the model reads $\mathcal{M}(v_b,
  \text{geometry}, \sigma_x, \sigma_y, t_r, v_p)$, whose parameters
  Section~\ref{sec:variables} varies.

  \subsection{Which Variables Matter}\label{sec:variables}

  The model of Section~\ref{sec:pwin} keeps $v_b$, $v_p$, $t_r$, and the spread
  $(\sigma_x, \sigma_y)$ as separate variables. To see which of them change the
  results, we measure each run by the highest $P_{\mathrm{win}}$ over the grid,
  with the grid, the $K$, and the common random numbers of
  Section~\ref{sec:pwin}. The base values are those of Section~\ref{sec:pwin},
  with $v_b = 12$~m/s for pickleball and a spread of $(0.5, 1.0)$ m. We first
  vary one variable at a time: $v_b$ from 5~m/s to the model's maximum ball
  speed $v_{\max}$ (73.14~m/s for tennis and 24.59~m/s for pickleball), $v_p$
  from 0 to 6~m/s, $t_r$ from 0 to 1 s, and each spread from 0.1 to 3 m. We then
  cross nine serves with the ends and the base values of $v_b$ and $v_p$, for
  both service sides and both sports. The upper ends of $v_p$ and of the spread
  bound the ranges of the model; they are not measured values. The experiments
  are in the accompanying code (Appendix~B).

  The top of the $v_b$ range is the $v_{\max}$ of the sport parameters of the
  accompanying code (Appendix~B), which is also the slider maximum of the
  notebooks, so every map of the sweep can be rebuilt there. For tennis it
  matches the record of Table~\ref{tab:court} to within 0.1~m/s; for pickleball
  it is below the 27.7~m/s record. The gap does not hide a drop: in this model a
  faster ball never lowers $P_{\mathrm{win}}$ (Section~\ref{sec:pwin}), so a
  sweep up to the record could only raise the highest $P_{\mathrm{win}}$ of
  pickleball, which is already 0.998 at 24.59~m/s.

  \textbf{The error spread matters most.} With a spread of 3 m the highest
  $P_{\mathrm{win}}$ falls to 0.827 for tennis and 0.695 for pickleball in
  $\sigma_x$, and to 0.955 and 0.727 in $\sigma_y$ (Fig.~\ref{fig:sweeps}). By
  contrast, $v_b$ and $v_p$ leave tennis at 1.000 and pickleball between 0.948
  and 0.998, and $t_r$ leaves pickleball between 0.996 and 0.998.

  \begin{figure}[htbp]
    \centering
    \includegraphics[width=0.62\textwidth]{variable_sensitivity_sweeps.png}
    \caption{Highest $P_{\mathrm{win}}$ when one variable moves and the others
      stay at their base values.}
    \label{fig:sweeps}
  \end{figure}

  \textbf{Player speed and ball speed matter mainly in one corner.} The highest
  $P_{\mathrm{win}}$ is at least 0.99 in 146 of the 162 tennis corner runs and
  in 108 of the 162 pickleball ones. The exception is a slow ball (5~m/s) with a
  fast server (6~m/s), where it falls to between 0.816 and 0.991 for tennis and
  between 0.047 and 0.705 for pickleball, depending on the serve. There the
  serve matters, and deep serves help the server (Fig.~\ref{fig:serve-depth}): a
  deep serve pushes the returner back, so the return flies farther and the
  server has more time to run. In pickleball the short serves leave 0.680 to
  0.705 and the deep ones 0.047 to 0.092, and the deep serve at the T is the
  lowest in both sports (on the tennis ad side, tied with the other deep
  serves). Outside that corner, pickleball still goes down to 0.904 at $v_p =
  6$~m/s with the default ball. The mean of $P_{\mathrm{win}}$ over the server's
  half moves more than its highest value, from 0.51 to 0.83 with $v_p$ in
  pickleball: the speeds change the cost of a poor aim more than the value of
  the best aim.

  \begin{figure}[htbp]
    \centering
    \includegraphics[width=0.62\textwidth]{variable_sensitivity_serve_depth.png}
    \caption{Highest $P_{\mathrm{win}}$ that the returner can get, by how deep
      the serve lands, for a slow ball and a fast server ($v_b = 5$~m/s, $v_p =
      6$~m/s). A lower value is better for the server.}
    \label{fig:serve-depth}
  \end{figure}

  \textbf{Reaction time changes little at realistic values.} Compared with $t_r
  = 0$, a reaction time of 0.2 s at $v_p = 1.5$~m/s, within the ranges of
  Section~\ref{sec:assumptions}, moves the best region by at most 0.26 m for
  tennis and 0.45 m for pickleball, less than a grid cell, over ball speeds from
  5~m/s to the fastest ball. The highest $P_{\mathrm{win}}$ does not move for
  tennis and moves by 0.82 percentage points for pickleball, and its mean over
  the half moves by at most 2.34 and 3.95 points. A reaction time matters more
  for a fast server, who loses the distance $v_p t_r$ while reacting: up to
  17.52 points on the highest $P_{\mathrm{win}}$ for pickleball at $v_p =
  6$~m/s. It is not an important variable at realistic values, and we keep it in
  the model for completeness. Without it the two speeds are not separate
  variables:

  \begin{prop}\label{prop:ratio}
    With $t_r = 0$ and a spread that does not depend on the ball speed,
    $P_{\mathrm{win}}$ of every target depends on $v_p$ and $v_b$ only through
    the ratio $v_p / v_b$.
  \end{prop}

  A proof sketch is in Appendix~A. In simulation, multiplying both speeds by the
  same factor gives identical maps, cell by cell, on 30 pairs of maps and on 500
  random settings. With $t_r > 0$ the same scaling changes the maps, since a
  second quantity enters, the reaction distance $v_p t_r$.
  Section~\ref{sec:speed} shows what changes when the spread grows with the
  ball speed.

  Two limits apply to these results. The highest $P_{\mathrm{win}}$ is 1.000 for
  many tennis settings, so it cannot rank them. And the spread does not depend
  on the ball speed here (Section~\ref{sec:pwin}), so a faster ball always looks
  better; Section~\ref{sec:speed} removes this limit.

  Details: research log, Experiments 3 and 4 (Appendix~B).

  \subsection{Speed Versus Accuracy}\label{sec:speed}

  With a spread that does not depend on the ball speed, a faster return never
  lowers $P_{\mathrm{win}}$, and the fastest ball always looks best. Players
  miss more when they hit harder: development players serving flat were less
  accurate at maximum power than at minimum power \cite{menayo2008serve}, and in
  the topspin forehands of college players the trade-off between ball speed and
  accuracy appeared in the hitting direction only \cite{sudo2024speedaccuracy}.
  The landing errors of advanced players also spread more along the sideline
  than across it \cite{delgadogarcia2019errors}, which supports a spread at a
  given speed that is larger in depth than in width.

  \gpnote{Sudo et al.\ (2024) correction}{%
    The publisher lists a correction to Sudo et al.\ (2024), published on
    2024-09-05 (doi
    \href{https://doi.org/10.1080/14763141.2024.2399923}{10.1080/14763141.2024.2399923}),
    which we could not read. Could you check whether it changes the result we
    cite twice in this section, that the trade-off between speed and accuracy
    appeared in the hitting direction only? \emph{Our recommendation:} keep
    the citation as it is, unless the correction touches that result.}

  We make the spread grow with the ball speed in the sweeps of this section
  only: the model still takes $v_b$, $\sigma_x$, and $\sigma_y$ as separate
  inputs, and each map sets the two spreads from its ball speed. With $u = (v_b
  - v_0) / v_0$, where the reference speed $v_0$ is the base ball speed of
  Section~\ref{sec:variables} (25~m/s for tennis and 12~m/s for pickleball),
  \begin{equation}
  \sigma_x = f(a_x) \, \sigma_{x,0}, \qquad \sigma_y = f(a_y) \, \sigma_{y,0},
  \end{equation}
  with the spread at $v_0$, $(\sigma_{x,0}, \sigma_{y,0}) = (0.5, 1.0)$ m, and
  one of four shapes of the factor: linear, $f(a) = 1 + a u$, the base; knee,
  $1 + a \max(0, u)$; parabola, $1 + a u^2$; and one-sided parabola, $1 + a
  \max(0, u)^2$. All four give the spread at $v_0$ for $v_b = v_0$, and $1 + a$
  times it for $v_b = 2 v_0$. They differ below $v_0$, where only the linear
  shape makes a slower ball more accurate. The depth gets $a_y = a$ and the
  width $a_x = a / 2$, our assumption between a width that does not change with
  the speed, as in \cite{sudo2024speedaccuracy}, and a width as sensitive as
  the depth. No study gives a value of $a$, so we sweep an assumed range, $a \in
  \{0, 0.25, 0.5, 0.75, 1\}$, with $v_b / v_0$ from 0.5 to 2 in steps of 0.1. We
  then change one value at a time, at $a = 0.5$ and 1: the other three shapes,
  $a_x = a_y$, the spread at $v_0$ halved and doubled, $v_p = 1.9$ and
  3.0~m/s, the deep serve at the T, and a second seed, at $a = 1$ only. The full
  run has 786 maps, with the grid, the $K$, and the common random numbers of
  Section~\ref{sec:pwin}, and the maps of a sweep share one seed. As for the
  best region, the best speed is a band: every speed whose highest
  $P_{\mathrm{win}}$ is within two standard errors of the highest over the
  sweep, marked at its middle. We call the best speed robust to a change if the
  two bands share at least one speed, a rule fixed before the runs.

  \gpnote{robustness test (Gdoc Q22)}{%
    Is ``the two bands share a speed'' the right robustness test? Under this
    rule no variant moves the band, but some bands share only an edge speed: in
    pickleball at $a = 1$ the other three shapes share only 0.8~$v_0$ with the
    linear base, while the marked speed moves from 0.7 to 1.0~$v_0$.
    \emph{Our recommendation:} keep the rule, since it was fixed before the
    runs, and report the shift of the marked speed next to it; a stricter rule
    would ask each marked speed to fall in the other band.}

  \textbf{Hitting much faster than $v_0$ does not pay.} In pickleball the best
  speed falls as $a$ grows (Fig.~\ref{fig:best-speed}): the band is 0.9 to
  2~$v_0$ at $a = 0$, where the fastest balls tie as in
  Section~\ref{sec:variables}, 0.5 to 1~$v_0$ at $a = 0.5$, and 0.5 to
  0.8~$v_0$ (6 to 9.6~m/s) at $a = 1$, marked at 0.7~$v_0$. The fastest ball,
  $2 v_0$, loses 2.9 points of the highest $P_{\mathrm{win}}$ at $a = 0.5$ and
  9.9 points at $a = 1$, and 2.7 to 9.9 points across the four shapes. In tennis
  the highest $P_{\mathrm{win}}$ is 1.000 at almost every speed, the saturation
  of Section~\ref{sec:variables}, so it cannot rank the speeds, but the mean
  over the server's half can. At $a = 1$ the mean is highest at 0.6~$v_0$
  (15~m/s), and the fastest ball loses 11.1 points of it; across the four shapes
  at $a \ge 0.5$ it loses 4.3 to 11.1 points.

  \textbf{The shape below $v_0$ decides the best pickleball speed.} With the
  linear shape the band reaches the slowest speed of the sweep, so the best
  pickleball speed is at or below $v_0 / 2$. With the knee and the two
  parabolas, which do not make a slower ball more accurate, it stays near
  $v_0$: at $a = 1$ their bands are 0.8 to 1.1 and 0.8 to 1.3~$v_0$, marked at
  1.0~$v_0$, against 0.7~$v_0$ for the linear shape. A slower ball gives the
  server more time, so it pays only if it is also more accurate. The bands still
  share one speed, 0.8~$v_0$, so the best speed is robust to the shape under our
  rule, but the marked speed moves by 0.3~$v_0$.

  \gpnote{accuracy below $v_0$ (Gdoc Q21)}{%
    Below $v_0$, does a slower ball become more accurate (the linear shape) or
    not (the knee)? The answer decides the best pickleball speed: at $a = 1$ it
    is marked at 0.7~$v_0$ with the linear shape and at 1.0~$v_0$ with the knee
    and the two parabolas, while above $v_0$ all four shapes agree that hitting
    much faster does not pay. You asked on Slack (2026-11-17) for a conditional
    claim; \#984 adds the threshold $a^\star$. \emph{Our recommendation:} keep
    the linear shape as the base, as decided on 2026-11-06, report the knee as
    the alternative, and ask the experts about play below $v_0$ with the
    pre-print.}

  \textbf{What the best speed is robust to.} The best speed is robust to $a_x =
  a_y$, to $v_p$, to the serve, and to the seed: they move the marked speed by
  at most 0.1~$v_0$, except $v_p = 3.0$~m/s in pickleball at $a = 0.5$, which
  moves it up by 0.3~$v_0$ while the bands share four speeds. The spread at
  $v_0$ matters more. Doubled, it pulls the band to the slow end, 0.5 to
  0.8~$v_0$ in tennis and only 0.5~$v_0$ in pickleball at $a = 1$, and the
  fastest ball loses 13.8 and 31.9 more points; halved, almost every speed ties.
  The shape below $v_0$ and the spread at $v_0$ are the main uncertainties of
  this section, and both call for data rather than more sweeps.

  \begin{figure}[htbp]
    \centering
    \includegraphics[width=0.62\textwidth]{speed_vs_accuracy_best_speed.png}
    \caption{Best ball speed $v_b / v_0$ against the coefficient $a$ of the
      spread: the linear base (line, with its band of tied speeds shaded) and
      the other three shapes at $a = 0.5$ and 1 (markers, with their bands as
      bars). In tennis the highest $P_{\mathrm{win}}$ is 1.000 at almost every
      speed, so the bands cover nearly the whole range; in pickleball the linear
      base falls to the slowest speed of the sweep, while the other shapes stay
      near $v_0$.}
    \label{fig:best-speed}
  \end{figure}

  \textbf{Proposition~1 with a spread that depends on the ball speed.} With $a >
  0$ the landing points depend on $v_b$ through $v_b / v_0$, so with $t_r = 0$,
  $P_{\mathrm{win}}$ depends on $v_p / v_b$ and on $v_b / v_0$ (Appendix~A).
  For the linear base at $t_r = 0$, doubling both speeds changes one cell by up
  to 0.23 in tennis and 0.32 in pickleball at $a = 1$, while doubling $v_0$ with
  them leaves every cell unchanged, at $a = 0$, 0.5, and 1 in both sports.

  Details: research log, Experiment 5 (Appendix~B).

  % ############################################################################
  \section{Reinforcement Learning Check of the Return}\label{sec:rl}

  Sections~\ref{sec:pwin} and~\ref{sec:variables} compute the best return from
  the model. This section asks whether an agent that only plays points, and
  sees only whether each return was a winner, finds the same answer. For
  rallies, or for a server who adapts, the space of situations is too large to
  map exhaustively, so learning is the practical route there; a setting whose
  answer is computed independently is where to check that the environment, the
  reward, and the training are right before learning goes beyond it. The agent
  reuses the code of the model, the sampling of a landing point and the winner
  event, but not its results: it never sees a map of $P_{\mathrm{win}}$, so
  agreement with Section~\ref{sec:pwin} is evidence that the learning setup
  works, not a consequence of how it was built.

  \subsection{Environment}\label{sec:env}

  We cast one return of serve as a Gymnasium environment
  \cite{towers2024gymnasium} in which the returner is the agent. One point is
  one shot, so an episode has a single step, and the problem is a contextual
  bandit \cite{lattimore2020bandit}: the serve is the context and the aim point
  is the arm.

  \begin{itemize}
    \item \textbf{Observation.} Six coordinates in meters, before rescaling
    (Section~\ref{sec:ppo}): the server $(x, y)$, the returner $(x, y)$, and the
    serve landing $(x_1, y_1)$. In this version the server stays at the center
    of the baseline and the returner hits from the serve landing, as in
    Section~\ref{sec:pwin}, so only the serve varies; the two players are in
    the observation so that later versions can move them.

    \item \textbf{Reset.} A serve drawn uniformly over the deuce service box.

    \item \textbf{Action.} One cell of the grid of Section~\ref{sec:pwin}, $10
    \times 15 = 150$ cells for tennis and $9 \times 10 = 90$ for pickleball. The
    return aims at the center of the cell.

    \item \textbf{Transition and reward.} One sample of the winner event $W_k$
    of Section~\ref{sec:pwin}, drawn from the random generator of the
    environment, with the reward $r = \mathbb{1}[W]$. The episode then ends.
  \end{itemize}

  The spread, the ball speed, $t_r$, and $v_p$ are fixed at the values of
  Section~\ref{sec:pwin}. Because the reward is the indicator of the winner
  event,
  \begin{equation}
  \EE[\, r | s, a \,] = P_{\mathrm{win}}(c_a | s) ,
  \end{equation}
  where $s$ is the serve and $c_a$ the cell of action $a$. The best action for
  a serve is therefore the best target of Section~\ref{sec:pwin}, and the
  actions within sampling error of it form its best region.

  \subsection{Reference Answers}\label{sec:reference}

  We compare the agent with two answers, each computed from scratch on the same
  grid and for the same serves. The first is a brute force through the
  environment: every action is played 1000 times, and the $k$-th play of every
  action uses the same reset seed, so the actions share common random numbers
  as in Section~\ref{sec:pwin}. The second is a fresh map of
  Section~\ref{sec:pwin}, with 5000 samples per cell and an independent seed.
  The environment reproduces Section~\ref{sec:pwin} if the two maps have a
  Spearman rank correlation of at least 0.95, and the best cell of each map lies
  in the best region of the other, up to ties: where the top cells of a map
  differ by less than their sampling error, its best cell is not well defined,
  and a miss among tied cells is not counted as a failure. The serves are 9 per
  sport, the centers of a $3 \times 3$ grid over the deuce box.

  \subsection{PPO Setup}\label{sec:ppo}

  The agent is PPO \cite{schulman2017ppo} as implemented in Stable-Baselines3
  \cite{raffin2021sb3}, version 2.9.2, with its default multilayer perceptron
  policy: separate policy and value networks, each with two hidden layers of 64
  tanh units, and a softmax over the actions. The hyperparameters are the
  defaults of the library: a learning rate of $3 \cdot 10^{-4}$, rollouts of
  2048 steps, 10 epochs per update on minibatches of 64, a clipping range of
  0.2, a value loss coefficient of 0.5, a gradient norm clipped at 0.5, and no
  entropy bonus. With one-step episodes the discount factor and the advantage
  estimator have no effect, and the advantage of an action is its reward minus
  the value of the serve.

  Two choices matter. First, the observation is rescaled to $[-1, 1]$ with the
  bounds of the court. The library feeds the raw observation to the network,
  and in meters the constant coordinate of the server, $y = -11.885$ m in
  tennis, saturated a quarter to a third of the first layer in early runs,
  which slowed learning. Second, each run trains for 75,776 steps, 37 rollouts,
  on random serves. On a reward that is almost always 1, PPO met the pass rule
  at about 75k steps in development runs and drifted away from it after about
  100k, so the length was fixed before the final runs. The agent searches the
  whole grid, with no action mask, no pruning with the results of
  Section~\ref{sec:pwin}, and no warm start. We train one agent per sport for
  each of three held-out seeds, not used during development.

  \subsection{Pass Rule}\label{sec:pass}

  The rule was written before the final runs, and a failing seed is reported as
  it is, with no retuning. Let $p^\star(s) = \max_a P_{\mathrm{win}}(c_a | s)$
  on the fresh map, and let $\pi(a | s)$ be the policy. We report:

  \begin{itemize}
    \item \textbf{Region.} Whether the greedy cell, $\argmax_a \pi(a | s)$, lies
    in the best region of the fresh map. The strict rule asks this for every
    serve, seed, and sport.

    \item \textbf{Expected regret.} The winner probability that the policy loses
    on average, in percentage points:
    \begin{equation}
    \mathrm{ER}(s) = \sum_a \pi(a | s) \,
    \big( p^\star(s) - P_{\mathrm{win}}(c_a | s) \big)
    \end{equation}

    \item \textbf{Greedy regret.} The same loss for the greedy cell alone.

    \item \textbf{Random baseline.} The expected regret of a uniformly random aim
    point, $p^\star(s)$ minus the mean of $P_{\mathrm{win}}$ over the cells.
  \end{itemize}

  The answer of Section~\ref{sec:pwin} is a region, not a point: many cells tie
  at the top, most of all in tennis. We therefore lead with the expected regret
  and with the region. A policy that spreads its probability over tied cells
  loses almost nothing, even when its single most likely cell changes from seed
  to seed.

  \subsection{Results}\label{sec:results}

  \textbf{The environment reproduces Section~\ref{sec:pwin}.} Over the 9 serves,
  the brute force and the fresh map have a Spearman rank correlation of 0.980
  to 0.982 for tennis and 0.990 to 0.994 for pickleball, and the best cell of
  each map lies in the best region of the other on 15 of 18 serves. The 3
  misses are the short pickleball serves, whose top cells tie within 0.2
  points, the sampling error of 1000 points per cell, so the criterion holds on
  all 18 serves up to ties (Fig.~\ref{fig:rl-maps}).

  \begin{figure}[htbp]
    \centering
    \includegraphics[width=\textwidth]{rl_env_vs_m2_pickleball.png}
    \caption{Pickleball, middle serve of the deuce box: P(winner) of every aim
      point from the brute force through the environment (left, 1000 points per
      cell, best region of 7 cells) and from a fresh Monte Carlo map of
      Section~\ref{sec:pwin} (right, 5000 samples per cell, best region of 4
      cells). Both put the best region, outlined in black, in the middle of the
      server's half behind the non-volley zone. Pickleball is the sharper test:
      its regions have 2 to 8 cells, against 40 to 45 of 150 in tennis. The
      tennis maps are in the research log (Appendix~B).}
    \label{fig:rl-maps}
  \end{figure}

  \textbf{PPO aims at the middle.} On all six runs, three seeds per sport, the
  expected regret is under 0.8 points for every serve, against 11 to 21 points
  for a random aim point. The greedy cell lies in the best region, in the middle
  of the server's half, on 43 of 54 serves. The strict rule passes on 4 of 6
  runs: seed 5 fails in both sports. In pickleball its greedy cell misses a
  region of 2 to 5 cells on 3 serves, by at most 0.38 points. In tennis the
  policy spreads its probability over the tied cells of the plateau, so its most
  likely cell is fragile: it misses the region on 8 of 9 serves, by up to 6.6
  points, while the expected regret of the same run stays under 0.8 points.
  With a 1-point tolerance on the greedy regret, 5 of 6 runs pass. The agent
  learns the answer of Section~\ref{sec:pwin} from the reward alone: the best
  return is the middle, not the lines.

  \subsection{What It Does Not Show Yet}\label{sec:notyet}

  \textbf{A weak test in tennis.} The tennis region holds 40 to 45 of 150
  cells, so a tennis pass is weak evidence; pickleball, with 2 to 8 cells, is
  the sharper test.

  \textbf{Drift.} On a reward that is almost always 1, PPO drifts after about
  100k steps. The fixed run length avoids the drift but does not explain it.

  \textbf{How much of the policy is in the region.} The greedy cell measures
  the region with a single point. The probability that the policy puts inside
  the best region would measure it directly; it needs a new training run, and
  we leave it to future work.

  \textbf{Scope.} The environment covers the deuce side, one shot, and a server
  who neither moves nor adapts. The ad side, rallies, self-play, and pruning
  the grid with the results of Section~\ref{sec:pwin} to speed up learning are
  future work (Section~\ref{sec:conclusion}).

  Details: research log, Experiment 6 (Appendix~B).

  % ############################################################################
  \section{Discussion and Limitations}\label{sec:discussion}

  \textbf{Neglected aerodynamics.} Measured drag and lift coefficients for
  tennis balls show that aerodynamic forces are not negligible relative to
  gravity, particularly at high speed and with topspin \cite{cross2014drag}.
  Because the no-drag model of Section~\ref{sec:trajectory} systematically
  shortens actual flight distance relative to the no-drag prediction for a
  given launch speed and angle, the trajectory-feasibility and flight-time
  estimates in this paper should be read as first-order approximations; the
  direction of the bias is known, but we have not quantified its magnitude for
  the specific speed and spin ranges of interest here.

  \textbf{Point-mass players and fixed speed.} Modeling players as points with a
  single maximum speed $v_p$ ignores acceleration and deceleration dynamics,
  the direction-dependence of movement speed (lateral versus forward recovery),
  and anticipation effects, whereby experienced players begin moving before
  ball contact based on early postural cues, which would effectively reduce
  $t_r$ below the simple-reaction-time baseline used here.

  \textbf{Error model coupling.} The stroke-error parameters $\sigma_\theta,
  \sigma_v, \sigma_\phi$ are treated in Section~\ref{sec:trajectory} as fixed,
  exogenous inputs, but in practice they likely increase with target
  difficulty, for example, wider angles or greater required pace probably carry
  larger execution error than safe, central targets. The current framework does
  not model this coupling; doing so would require calibrating $\sigma(\cdot)$
  as a function of target geometry from empirical stroke data, which is a
  natural extension; Section~\ref{sec:speed} models the pace part of it in the
  2D model, with an assumed coefficient. The 2D model has the same limit: its
  spread $(\sigma_x, \sigma_y)$ does not depend on how far the returner runs
  before hitting, a coupling that would make the serve a strategic choice.

  \textbf{Straight-line flight and no bounce.} The 2D model of
  Section~\ref{sec:pwin} flies the ball in a straight line at a constant ground
  speed, and asks the server to reach the landing point by the time the ball
  lands. Real returns slow down and bounce, and the server can play the ball
  after the bounce, which gives them more time. The model therefore understates
  the server's reach, and with it the value of aiming away from the server: the
  main result holds as long as that reach stays small compared with the spread
  of the return.

  \textbf{Saturation.} At the default settings many aim points reach a
  $P_{\mathrm{win}}$ close to 1 in tennis, so the highest value cannot rank
  settings (Section~\ref{sec:variables}), the best region is large
  (Fig.~\ref{fig:winner-map}), and a learned policy can pass in tennis without
  being precise (Section~\ref{sec:rl}). A faster server, a slower ball, or a
  larger spread sharpens the comparison; whether the defaults are realistic or
  need calibrating is open.

  \textbf{Deterministic reachability.} Section~\ref{sec:reach} treats $R(c)$ as
  a strict indicator. In practice a player's reaction time and top speed are
  themselves variable, so a graded, probabilistic reachability model, for
  example, a smooth function of the margin $T_f(c) - t_r - d(O,c)/v_p$ rather
  than a step function, would likely produce a more realistic composite score
  near the boundary of the reach envelope.

  \textbf{Grid resolution versus variance.} A finer grid better approximates the
  continuous placement decision but requires more Monte Carlo samples per cell
  to control estimation variance in $P_{\mathrm{in}}$, and the resulting larger
  payoff matrix increases the cost of the game-theoretic solution step in
  Section~\ref{sec:game}. The 2D model uses common random numbers and a fixed
  grid ($10 \times 15$ cells for tennis and $9 \times 10$ for pickleball,
  Section~\ref{sec:pwin}), and we do not study this trade-off quantitatively
  here.

  % ############################################################################
  \section{Conclusion and Future Work}\label{sec:conclusion}

  We proposed a reduced-order framework for optimal shot placement in tennis and
  pickleball and computed its two-dimensional instance for the return of serve.
  The winner probability of every aim point is estimated by Monte Carlo as one
  joint event, on common random numbers, and the best placement is a region of
  statistically tied aim points. At the default settings the best return is the
  middle of the server's half in both sports, not a spot near the lines:
  landing in matters more than landing far from the opponent. The error spread
  matters most; the ball and player speeds matter mainly with a slow ball and a
  fast server, and with no reaction time they act only through their ratio
  (Proposition~1). When a faster return is less accurate, hitting much faster
  than the default ball speed does not pay, while the best slower speed depends
  on how accurate a slower ball is (Section~\ref{sec:speed}). A one-step
  reinforcement learning environment built on the same model reproduces the
  winner probability map from the reward alone, and a PPO agent trained on it
  aims at the same middle region. The general framework, a trajectory model, a
  grid score with reachability, and a zero-sum game connected to empirical
  evidence of minimax play in professional tennis \cite{walker2001minimax},
  frames how this model extends to full rallies, and the same parametrized
  model covers both sports.

  The results are computational: the parameters come from literature ranges and
  round values, not from tracking data, so they show how the model behaves
  rather than how players play. The most immediate next steps are: (i)
  measuring how the spread of a return grows with the ball speed, which
  Section~\ref{sec:speed} sweeps with an assumed coefficient, and making it grow
  with the distance the returner runs; (ii) calibrating the spread $(\sigma_x,
  \sigma_y)$ and the reachability parameters $(t_r, v_p)$ from tracking data or
  from experts, rather than from literature ranges; (iii) replacing the
  deterministic reachability with a graded, probabilistic version; (iv)
  extending the reinforcement learning environment to the ad side, to rallies
  in which both players move, and to self-play, which leads to the post-shot
  positioning game of Section~\ref{sec:postshot}; (v) explaining the simulated
  maps in closed form, for example by symbolic regression; and (vi) extending
  the geometric model to doubles court dimensions and validating the
  minimax-equilibrium hypothesis of Section~\ref{sec:wimbledon} directly on
  placement data, rather than only on serve-direction data.

  % ############################################################################
  \section*{Appendix~A: Proof Sketch of Proposition~1}

  Take the model of Section~\ref{sec:pwin}. The landing points $(X_k, Y_k)$ are
  drawn from the target and the spread, so their distribution does not depend
  on $v_b$ or $v_p$, and whether a sample lands in the server's half $H$
  depends only on $(X_k, Y_k)$. With $d_k = \lVert (X_k, Y_k) - (x_1, y_1)
  \rVert$ the flight distance and $r = v_p / v_b$, the reach rule of
  Section~\ref{sec:pwin} reads
  \begin{equation}
  \begin{aligned}
  \lVert (X_k, Y_k) - P \rVert
  &\le v_p \max\left(0, \, \frac{d_k}{v_b} - t_r\right) \\
  &= \max(0, \, r \, d_k - v_p t_r) .
  \end{aligned}
  \end{equation}
  So each event $W_k$ depends on the two speeds only through $r$ and the
  reaction distance $v_p t_r$. With $t_r = 0$ the reaction distance is 0 and
  only $r$ is left. $P_{\mathrm{win}}$ is the probability of this event, so it
  depends on $v_p$ and $v_b$ only through $r$, and so do the best region and
  the best target, which are functions of the map. With $t_r > 0$ the map
  depends on both $r$ and $v_p t_r$.

  With the spread of Section~\ref{sec:speed} the first step changes: the
  spread, and with it the distribution of $(X_k, Y_k)$, depends on $v_b$ through
  $v_b / v_0$. The reach rule does not change, so each event $W_k$ depends on
  the speeds through $r$, $v_p t_r$, and $v_b / v_0$, and with $t_r = 0$ through
  $r$ and $v_b / v_0$. Multiplying $v_p$, $v_b$, and $v_0$ by the same factor
  leaves all three unchanged, so it leaves the map unchanged, for every shape of
  the factor and any $a_x$ and $a_y$.

  % ############################################################################
  \section*{Appendix~B: Experiment Index}

  The paper holds the model and the reasoning; the runs, the figures, the
  settings strings, and the decisions behind each result are in the project's
  research log, a Google Doc whose Research results tab has one entry per
  experiment:

  \begin{sloppypar}
    \noindent\url{https://docs.google.com/document/d/19KGM-1a9U5E7CrE4ltljYz7IVtMfNuRzrqgGkkuXUSQ/edit?tab=t.0}
  \end{sloppypar}

  The accompanying code holds:

  \begin{itemize}
    \item \path{notebooks/racket_sports_NN_*.ipynb}: notebook $NN$ of the table
    below

    \item \path{notebooks/racket_sports_01_shot_probability_utils.py}: the model
    of Sections~\ref{sec:pin} and~\ref{sec:pwin}

    \item \path{notebooks/racket_sports_06_rl_return_utils.py}: the environment
    of Section~\ref{sec:rl}

    \item \path{notebooks/racket_sports_07_speed_vs_accuracy_utils.py}: the
    spread model and the sweeps of Section~\ref{sec:speed}

    \item \path{racket_params.py}: the sport parameters, including $v_{\max}$

    \item \path{figures/make_code_figures.py}, next to the paper:
    Fig.~\ref{fig:winner-map} and Fig.~\ref{fig:rl-maps};
    \path{figures/make_figures.py} draws Fig.~\ref{fig:sweeps} and
    Fig.~\ref{fig:serve-depth} from the data that notebook 05 writes, and
    Fig.~\ref{fig:best-speed} from the data that notebook 07 writes
  \end{itemize}

  \begin{table}[htbp]
    \centering
    \caption{Where the details of each result are kept}
    \label{tab:index}
    \begin{tabular}{llll}
      \toprule
      Section & Result & Notebook & Research log \\
      \midrule
      \ref{sec:pin} & In-bounds probability & 01, Part 1
        & \href{https://docs.google.com/document/d/19KGM-1a9U5E7CrE4ltljYz7IVtMfNuRzrqgGkkuXUSQ/edit?tab=t.0}{Experiment 1} \\
      \ref{sec:pwin} & Winner probability map & 01, Part 2
        & \href{https://docs.google.com/document/d/19KGM-1a9U5E7CrE4ltljYz7IVtMfNuRzrqgGkkuXUSQ/edit?tab=t.0}{Experiment 2} \\
      \ref{sec:variables} & Sweeps and serves & 05
        & \href{https://docs.google.com/document/d/19KGM-1a9U5E7CrE4ltljYz7IVtMfNuRzrqgGkkuXUSQ/edit?tab=t.0}{Experiment 3} \\
      \ref{sec:variables} & Speed ratio, reaction time & 05, Part 3
        & \href{https://docs.google.com/document/d/19KGM-1a9U5E7CrE4ltljYz7IVtMfNuRzrqgGkkuXUSQ/edit?tab=t.0\#heading=h.xvdn4n7zo860}{Experiment 4} \\
      \ref{sec:speed} & Speed vs accuracy & 07
        & \href{https://docs.google.com/document/d/19KGM-1a9U5E7CrE4ltljYz7IVtMfNuRzrqgGkkuXUSQ/edit?tab=t.0\#heading=h.qrto9so35bnx}{Experiment 5} \\
      \ref{sec:rl} & Reinforcement learning & 06
        & \href{https://docs.google.com/document/d/19KGM-1a9U5E7CrE4ltljYz7IVtMfNuRzrqgGkkuXUSQ/edit?tab=t.0\#heading=h.kntokyu5w697}{Experiment 6} \\
      \bottomrule
    \end{tabular}
  \end{table}

  \bibliographystyle{unsrtnat}
  \bibliography{references}
\end{document}
